% ================ CHAPITRE 5 — VOLET DÉCISIONNEL =====================
\chapter{Volet décisionnel}
\label{chap:decisionnel}

\section*{Introduction}
Ce chapitre présente le \textbf{volet décisionnel} de la solution. Il met en place
une chaîne \textbf{ETL} qui consolide les données opérationnelles dans un
\textbf{datamart}, restitue les indicateurs au moyen d'un tableau de bord embarqué
et de rapports Power BI, puis introduit une brique d'\textbf{intelligence
artificielle} : la segmentation des agences par apprentissage non supervisé.

% ----------------------------------------------------------------------
\section{Chaîne décisionnelle et processus ETL}
Les données nécessaires au pilotage sont dispersées dans les tables opérationnelles
et exprimées à des grains hétérogènes. Le processus \textbf{ETL} (\emph{Extract --
Transform -- Load}) les consolide en un datamart agrégé par agence. La
figure~\ref{fig:etl_chaine} présente cette chaîne décisionnelle.

\begin{figure}[htbp]
  \centering
  \includegraphics[width=\textwidth]{etl_chaine.png}
  \caption{Chaîne décisionnelle : du système opérationnel au datamart et à sa restitution}
  \label{fig:etl_chaine}
\end{figure}

Trois étapes s'enchaînent : \textbf{Extract} (lecture des sources
\texttt{AGENCE}, \texttt{B\_UTILISATEURS}, \texttt{CLIENT\_BTK}, \texttt{B\_OBJECTIF},
\texttt{POINTAGE}), \textbf{Transform} (nettoyage, typage et agrégation par agence)
et \textbf{Load} (chargement du datamart en étoile). Ce datamart constitue une
\textbf{source unique consolidée}, partagée par la restitution (tableau de bord et
Power BI) et par la segmentation des agences présentée à la
section~\ref{sec:segmentation}.

% ----------------------------------------------------------------------
\section{Modèle en étoile de l'entrepôt}
L'entrepôt est organisé selon une logique en \textbf{étoile}
(figure~\ref{fig:bi}) : autour de tables de faits (production, effectifs, clients,
présence), des dimensions descriptives (agence, gestionnaire, période, type de
client) permettent l'analyse multidimensionnelle. Les vues \texttt{V\_BI\_*}
matérialisent cette couche sémantique pour l'outil de restitution.

\begin{figure}[htbp]
  \centering
  \includegraphics[width=0.95\textwidth]{bi.png}
  \caption{Modèle décisionnel en étoile (vues \texttt{V\_BI\_*})}
  \label{fig:bi}
\end{figure}

Le tableau~\ref{tab:kpi} récapitule les principaux indicateurs (KPI) retenus.

\begin{table}[htbp]
\centering
\caption{Indicateurs de performance (KPI) du volet décisionnel}
\label{tab:kpi}
\renewcommand{\arraystretch}{1.25}
\begin{tabularx}{\textwidth}{>{\bfseries}p{4.2cm}Xp{2.6cm}}
\toprule
\textmd{\textbf{Indicateur}} & \textbf{Description} & \textbf{Source}\\
\midrule
Effectifs par agence & Nombre d'employés par agence & \texttt{V\_BI\_EMPLOYES}\\
Répartition des rôles & Gestionnaires / conseillers / hors commercial & \texttt{V\_BI\_EMPLOYES}\\
Portefeuille clients & Clients par type et par agence & \texttt{V\_BI\_CLIENTS}\\
Ouvertures de comptes & Comptes chèques + épargne + courants & \texttt{V\_BI\_OBJECTIFS}\\
Production de crédits & Conso, immobilier, investissement & \texttt{V\_BI\_OBJECTIFS}\\
Collecte d'épargne & Épargne additionnelle collectée & \texttt{V\_BI\_OBJECTIFS}\\
Taux de présence & Présents, retards et absences par agence et période & \texttt{V\_BI\_POINTAGE}\\
\bottomrule
\end{tabularx}
\end{table}

% ----------------------------------------------------------------------
\section{Mise en œuvre de la chaîne ETL en Python}

\subsection{Export des données sources}
La chaîne ETL peut se connecter directement à Oracle ou, à défaut, consommer des
fichiers CSV. Pour ce second mode, un script \textbf{SQLcl}
(\texttt{etl/export\_sources.sql}) exporte les cinq tables sources vers des CSV
horodatés (extrait~\ref{lst:export}). Chaque \texttt{SPOOL} redirige le résultat
d'une requête vers un fichier, placé ensuite dans \texttt{etl/source/}.

\begin{lstlisting}[style=pycode,language=SQL,caption={Export des sources vers CSV (SQLcl)},label={lst:export}]
SET SQLFORMAT CSV
SPOOL agences.csv
SELECT SK_AGENCE, LIBELLE_AGENCE, DISTRICT FROM AGENCE;
SPOOL OFF

SPOOL objectifs.csv
SELECT SK_AGENCE,
       NVL(SOUSCRIPTION_COMPTE_CHEQUES_OA,0)
         + NVL(SOUSCRIPTION_COMPTE_EPARGNES_OA,0)
         + NVL(SOUSCRIPTION_COMPTE_COURANTS_OA,0) AS COMPTES,
       NVL(PRODUCTION_CREDITS_CONSO_OA,0)
         + NVL(PRODUCTION_CREDITS_IMMO_OA,0)
         + NVL(PRODUCTION_CREDITS_INVESTISSEMENT_OA,0) AS CREDITS,
       NVL(EPARGNE_ADD_OA,0) AS EPARGNE
FROM   B_OBJECTIF;
SPOOL OFF
\end{lstlisting}

\subsection{Agrégation et chargement}
La chaîne ETL est implémentée en \textbf{Python} (script \texttt{etl\_agences.py})
à l'aide de la bibliothèque \textbf{pandas}~\cite{pandas}. Elle est organisée en
trois fonctions : \texttt{extract()} lit les données sources (connexion Oracle via
le connecteur \texttt{oracledb} ou, à défaut, export CSV) ; \texttt{transform()}
nettoie et \textbf{agrège les données par agence} pour produire les sept
indicateurs de pilotage ; \texttt{load()} écrit le datamart en étoile
(\texttt{dim\_agence}, \texttt{fait\_agence}) et le fichier réutilisé par la
segmentation. L'extrait~\ref{lst:etl_transform} illustre le cœur de l'agrégation.

\begin{lstlisting}[style=pycode,caption={Agrégation des indicateurs par agence (étape Transform)},label={lst:etl_transform}]
# effectif et nombre de gestionnaires par agence
eff = employes.groupby("SK_AGENCE").agg(
    effectif=("SK_UTILISATEUR", "count"),
    nb_gestionnaires=("EST_GESTIONNAIRE", "sum")).reset_index()
# portefeuille clients
cli = clients.groupby("SK_AGENCE").size().reset_index(name="nb_clients")
# production (comptes, credits, epargne)
obj = objectifs.groupby("SK_AGENCE").agg(
    total_comptes=("COMPTES", "sum"),
    production_credits=("CREDITS", "sum"),
    collecte_epargne=("EPARGNE", "sum")).reset_index()
# taux de presence issu du pointage
poi = pointages.assign(
        present=pointages["STATUT"].isin(["PRESENT", "RETARD"]).astype(int)) \
    .groupby("SK_AGENCE").agg(taux_presence=("present", "mean")).reset_index()
\end{lstlisting}

L'exécution du script produit le datamart sous \texttt{etl/entrepot/}
(\texttt{dim\_agence.csv}, \texttt{fait\_agence.csv}) ainsi que le fichier
consolidé \texttt{clustering/data/agences.csv}. Le pipeline est \textbf{rejouable}
et journalise le nombre d'enregistrements traités à chaque étape.

% ----------------------------------------------------------------------
\section{Tableau de bord et restitution Power BI}
Le tableau de bord embarqué (\texttt{home.jsp}, alimenté par le
\texttt{HomeServlet}) comporte une rangée de \textbf{cartes d'indicateurs}
(employés, clients, présents du jour, taux de présence, demandes en attente) puis
des \textbf{graphiques} ApexCharts (effectif par agence, répartition des rôles,
présence du jour). En complément, les vues \texttt{V\_BI\_*} sont connectées à
\textbf{Power BI} pour des rapports interactifs. La figure~\ref{fig:ui_dashboard}
présente le tableau de bord embarqué.

\begin{figure}[htbp]
  \centering
  \includegraphics[width=\textwidth]{screens/dashboard.png}
  \caption{Tableau de bord décisionnel embarqué de l'application}
  \label{fig:ui_dashboard}
\end{figure}

% ----------------------------------------------------------------------
\section{Rapport décisionnel sous Power BI}
Au-delà du tableau de bord embarqué, un \textbf{rapport décisionnel dédié} a été
développé sous \textbf{Microsoft Power BI}, connecté aux données consolidées. Il
est organisé en \textbf{quatre pages} analytiques thématiques --- \emph{Résumé},
\emph{Clients}, \emph{Commercial} et \emph{Présence} --- reliées par une barre de
navigation, chacune munie de \textbf{segments} de filtrage (district, période,
type de client) qui recalculent l'ensemble des visuels d'un simple clic. La
\textbf{page de synthèse} (\emph{Résumé}) regroupe les indicateurs clés du réseau
(clients, production de crédits, taux de présence, ouvertures de comptes), la
répartition des rôles et des présences, l'évolution mensuelle de la production et
un classement des agences. Les trois pages suivantes approfondissent
respectivement l'analyse des clients, la performance commerciale et l'assiduité.

\subsection{Indicateurs et mesures (DAX)}
Les indicateurs restitués sont calculés au moyen de \textbf{mesures DAX} définies
sur le modèle de données. Le tableau~\ref{tab:dax} en récapitule les principales.

\begin{table}[htbp]
\centering
\caption{Principales mesures DAX du rapport Power BI}
\label{tab:dax}
\renewcommand{\arraystretch}{1.25}
\begin{tabularx}{\textwidth}{>{\bfseries}p{3.8cm}X}
\toprule
\textmd{\textbf{Mesure}} & \textbf{Définition}\\
\midrule
Nb Employés & Nombre distinct d'employés (\texttt{SK\_UTILISATEUR}).\\
Nb Clients & Nombre distinct de clients (\texttt{SK\_CLIENT}).\\
Total Comptes & Somme des ouvertures de comptes (chèques + épargne + courants).\\
Total Crédits & Somme de la production de crédits (conso + immobilier + investissement).\\
Total Épargne & Épargne additionnelle collectée.\\
Total Packs \& Cartes & Souscriptions de packs (particulier, pro) et de cartes.\\
Taux Présence & Part des pointages au statut « présent » ou « retard ».\\
Nb Présences / Retards / Absences & Nombre de pointages par statut.\\
Rôle \emph{(colonne calculée)} & Gestionnaire / Conseiller / Hors commercial, selon \texttt{EST\_GESTIONNAIRE} et l'agence.\\
\bottomrule
\end{tabularx}
\end{table}

\subsection{Analyse des clients}
La page \emph{Clients} (figure~\ref{fig:pbi_clients}) segmente le portefeuille selon
plusieurs axes : le \textbf{type de client}, le \textbf{sexe}, la \textbf{région}
et le \textbf{gouvernorat} (histogramme des principaux gouvernorats), complétés par
un classement des agences selon leur nombre de clients. Elle offre une lecture
\textbf{géographique et démographique} du portefeuille, absente du tableau de bord
opérationnel.

\begin{figure}[htbp]
  \centering
  \includegraphics[width=\textwidth]{screens/pbi_clients.png}
  \caption{Rapport Power BI --- page « Clients » (segmentation et territoire)}
  \label{fig:pbi_clients}
\end{figure}

\subsection{Performance commerciale}
La page \emph{Commercial} (figure~\ref{fig:pbi_commercial}) détaille les
\textbf{objectifs} : les ouvertures de comptes par type, la production de crédits
par nature, l'effort d'équipement (EER), l'\textbf{évolution mensuelle} de la
production et un classement des agences par production. Elle constitue l'outil de
\textbf{pilotage commercial} du réseau.

\begin{figure}[htbp]
  \centering
  \includegraphics[width=\textwidth]{screens/pbi_commercial.png}
  \caption{Rapport Power BI --- page « Commercial » (objectifs et production)}
  \label{fig:pbi_commercial}
\end{figure}

\subsection{Suivi de la présence}
La page \emph{Présence} (figure~\ref{fig:pbi_presence}) suit l'\textbf{assiduité}
des employés : le taux de présence, la répartition des statuts (présent, retard,
absent) et de la source de pointage (automatique, manuel), l'\textbf{évolution
mensuelle} du taux de présence et la liste des \textbf{agences à surveiller}
(présence la plus faible). Elle prolonge le module de pointage par une vue
\textbf{ressources humaines} consolidée.

\begin{figure}[htbp]
  \centering
  \includegraphics[width=\textwidth]{screens/pbi_presence.png}
  \caption{Rapport Power BI --- page « Présence » (assiduité et suivi RH)}
  \label{fig:pbi_presence}
\end{figure}

\subsection{Navigation et interactivité}
Les quatre pages sont reliées par un \textbf{navigateur} qui joue le rôle de menu.
Chaque page est équipée de \textbf{segments} (district, période, type de client)
et d'un bouton \textbf{RÉINITIALISER} : un clic sur une valeur recalcule
instantanément l'ensemble des visuels de la page (\emph{filtrage croisé}), et la
sélection d'une part d'un graphique filtre à son tour les autres. Cette
\textbf{interactivité} transforme le rapport en un véritable outil d'exploration :
le décideur peut, par exemple, isoler un district puis observer simultanément le
portefeuille, la production et l'assiduité correspondants. Un \textbf{thème
sombre} homogène assure la lisibilité et la cohérence visuelle de l'ensemble.

% ----------------------------------------------------------------------
\section{Segmentation des agences par apprentissage non supervisé}
\label{sec:segmentation}
Cette dernière section introduit une brique d'\textbf{intelligence artificielle} :
la segmentation des agences par \textbf{apprentissage non supervisé}
(\emph{clustering}). À partir du datamart produit par la chaîne ETL, elle regroupe
automatiquement les agences en profils homogènes afin d'appuyer le ciblage,
l'allocation des ressources et la définition d'objectifs différenciés.

\subsection{Données et variables}
Les données ne sont pas extraites directement de la base opérationnelle : elles
proviennent du \textbf{datamart} produit par la chaîne ETL (table de faits
\texttt{fait\_agence}). Chaque agence y est décrite par sept indicateurs agrégés
(tableau~\ref{tab:variables_clustering}), ce qui garantit la cohérence entre le
tableau de bord, Power BI et l'analyse non supervisée.

\medskip\noindent\emph{Remarque.} En l'absence d'un export des données de
production, la segmentation présentée ici est \textbf{illustrée sur un jeu de
données représentatif} du réseau, généré par la chaîne ETL pour reproduire la
structure et les ordres de grandeur des agences ; le pipeline bascule
automatiquement sur les données réelles dès qu'un export est fourni
(\texttt{etl/source/}). Les résultats ci-dessous illustrent donc la
\emph{démarche} et restent reproductibles sur les données réelles.

\begin{table}[htbp]
\centering
\caption{Variables de segmentation (par agence)}
\label{tab:variables_clustering}
\renewcommand{\arraystretch}{1.2}
\begin{tabularx}{\textwidth}{>{\bfseries}p{4cm}X}
\toprule
\textmd{\textbf{Variable}} & \textbf{Description}\\
\midrule
effectif & Nombre d'employés de l'agence\\
nb\_gestionnaires & Nombre de gestionnaires\\
nb\_clients & Taille du portefeuille clients\\
total\_comptes & Total des ouvertures de comptes\\
production\_credits & Production de crédits\\
collecte\_epargne & Collecte d'épargne\\
taux\_presence & Taux de présence (issu du pointage)\\
\bottomrule
\end{tabularx}
\end{table}

\subsection{Prétraitement}
Les variables étant exprimées dans des ordres de grandeur très différents (un
effectif de quelques dizaines face à des milliers de comptes), une
\textbf{standardisation} (centrage-réduction, \texttt{StandardScaler}) est
appliquée. Sans elle, les variables de forte amplitude domineraient le calcul des
distances et fausseraient le regroupement.

\subsection{Détermination du nombre de clusters}
Le nombre de clusters $k$ étant inconnu a priori, il est déterminé en combinant la
\textbf{méthode du coude} (inertie intra-cluster) et le \textbf{score de
silhouette} (figure~\ref{fig:choix_k}).

\begin{figure}[htbp]
  \centering
  \includegraphics[width=\textwidth]{clustering_choix_k.png}
  \caption{Choix du nombre de clusters : méthode du coude et score de silhouette}
  \label{fig:choix_k}
\end{figure}

Les deux indicateurs concordent : l'inertie présente un \textbf{coude marqué à
$k=3$} et le score de silhouette y atteint son \textbf{maximum}
($\approx 0{,}74$). La segmentation en \textbf{trois clusters} est donc retenue.

\subsection{Comparaison des modèles de clustering}
Quatre algorithmes ont été comparés : \textbf{K-Means}, \textbf{classification
ascendante hiérarchique} (agglomératif), \textbf{modèle de mélange gaussien} (GMM)
et \textbf{DBSCAN}, à l'aide de trois indices internes
(tableau~\ref{tab:comparaison_clustering}, figure~\ref{fig:comparaison_clustering}).

\begin{table}[htbp]
\centering
\caption{Comparaison des modèles de clustering}
\label{tab:comparaison_clustering}
\renewcommand{\arraystretch}{1.25}
\begin{tabularx}{\textwidth}{>{\bfseries}lcccc}
\toprule
\textmd{\textbf{Modèle}} & \textbf{Nb clusters} & \textbf{Silhouette $\uparrow$} &
\textbf{Davies-Bouldin $\downarrow$} & \textbf{Calinski-Harabasz $\uparrow$}\\
\midrule
K-Means        & 3 & \textbf{0,736} & \textbf{0,350} & \textbf{391}\\
Agglomératif   & 3 & 0,736 & 0,350 & 391\\
GMM            & 3 & 0,736 & 0,350 & 391\\
DBSCAN         & 3 & 0,736 & 0,350 & 391\\
\bottomrule
\end{tabularx}
\end{table}

\begin{figure}[htbp]
  \centering
  \includegraphics[width=0.7\textwidth]{clustering_comparaison.png}
  \caption{Comparaison des modèles selon le score de silhouette}
  \label{fig:comparaison_clustering}
\end{figure}

Les quatre algorithmes \textbf{convergent vers la même segmentation} : ils
identifient les trois mêmes clusters et obtiennent des indices identiques
(silhouette $0{,}736$, Davies-Bouldin $0{,}350$, Calinski-Harabasz $391$). Cette
concordance, y compris pour des approches de nature différente (partitionnement,
hiérarchique, probabiliste et à base de densité), confirme la \textbf{robustesse}
des trois segments. Le \textbf{K-Means} est retenu : il atteint le meilleur score
tout en restant simple et interprétable.

\subsection{Résultats et interprétation}
La figure~\ref{fig:clusters_pca} projette les agences dans le plan des deux
premières composantes principales (PCA), colorées selon leur cluster ; la première
composante, qui résume la \emph{taille} de l'agence, explique l'essentiel de la
variance. Les trois segments sont nettement séparés.

\begin{figure}[htbp]
  \centering
  \includegraphics[width=0.78\textwidth]{clustering_pca.png}
  \caption{Segmentation des agences (K-Means) projetée par PCA}
  \label{fig:clusters_pca}
\end{figure}

Le profil moyen de chaque segment (tableau~\ref{tab:profils_clusters}) permet de
les caractériser.

\begin{table}[htbp]
\centering
\caption{Profil moyen des segments d'agences}
\label{tab:profils_clusters}
\renewcommand{\arraystretch}{1.25}
\begin{tabularx}{\textwidth}{>{\bfseries}Xcccccc}
\toprule
\textmd{\textbf{Segment}} & \textbf{Nb} & \textbf{Effectif} & \textbf{Clients} &
\textbf{Comptes} & \textbf{Prod. crédits} & \textbf{Présence}\\
\midrule
Grandes agences   & 7  & 116 & 908 & 2\,333 & 1\,734 & 0,98\\
Moyennes agences  & 16 & 54  & 438 & 1\,076 & 823  & 0,94\\
Petites agences   & 22 & 21  & 160 & 431  & 321  & 0,87\\
\bottomrule
\end{tabularx}
\end{table}

On distingue ainsi un segment de \textbf{grandes agences} (pôles majeurs à fort
effectif et forte production), un segment d'\textbf{agences moyennes} (activité
intermédiaire) et un segment de \textbf{petites agences} (effectif et activité
réduits, taux de présence plus faible).

\subsection{Apport décisionnel}
Cette segmentation constitue un outil d'\textbf{aide à la décision} : elle permet
d'adapter les actions à chaque profil. Le tableau~\ref{tab:reco} décline des
recommandations de pilotage par segment.

\begin{table}[htbp]
\centering
\caption{Recommandations de pilotage par segment}
\label{tab:reco}
\renewcommand{\arraystretch}{1.25}
\begin{tabularx}{\textwidth}{>{\bfseries}p{3.4cm}X}
\toprule
\textmd{\textbf{Segment}} & \textbf{Recommandations}\\
\midrule
Grandes agences & Capitaliser sur la performance ; fixer des objectifs ambitieux ;
en faire des agences de référence et diffuser leurs bonnes pratiques.\\
Moyennes agences & Fort potentiel de croissance ; renforcer l'effort d'équipement
(EER) et la collecte ; assigner des objectifs progressifs.\\
Petites agences & Accompagnement ciblé ; suivi rapproché de l'assiduité (taux de
présence plus faible) ; mutualisation des ressources et montée en compétence.\\
\bottomrule
\end{tabularx}
\end{table}

En résumé, la segmentation permet d'\textbf{allouer les ressources} selon la
taille, de \textbf{différencier les objectifs} et d'\textbf{accompagner
spécifiquement} les agences les plus fragiles. Elle prolonge le volet décisionnel
et illustre l'apport de l'\textbf{intelligence artificielle} au pilotage du réseau.

\section*{Conclusion}
Ce chapitre a doté la solution d'un volet décisionnel complet : une chaîne ETL
alimente un datamart consolidé, restitué par le tableau de bord et par Power BI,
et exploité par une segmentation des agences qui fait émerger trois profils
cohérents et exploitables pour la décision.
